\subsection{مبانی نظری و فرمول‌بندی ریاضی}
فرض کنید $\mathbf{X} \in \mathbb{R}^{B \times N \times D}$ نمایانگر نمایش‌های پنهان ورودی برای دسته به اندازه $B$، طول توالی $N$ و بعد پنهان $D$ باشد. تانسورهای پرس‌وجو ($\mathbf{Q}$)، کلید ($\mathbf{K}$) و مقدار ($\mathbf{V}$) با تعداد سرهای $H$ و بعد هر سر $d$ ($D = H \cdot d$) تعریف می‌شوند:
\begin{equation}
\mathbf{Q}_h = \mathbf{X} \mathbf{W}_h^Q, \quad \mathbf{K}_h = \mathbf{X} \mathbf{W}_h^K, \quad \mathbf{V}_h = \mathbf{X} \mathbf{W}_h^V
\end{equation}

در این ماژول، دو شاخه محاسباتی موازی جریان می‌یابند:
\begin{enumerate}
    \item \textbf{شاخه توجه استاندارد (\lr{Standard Softmax Branch}):}
    \begin{equation}
    \mathbf{O}_h^{\text{softmax}} = \operatorname{Softmax}\left(\frac{\mathbf{Q}_h \mathbf{K}_h^T}{\sqrt{d}} + \mathbf{M}\right) \mathbf{V}_h
    \end{equation}
    که در آن $\mathbf{M}$ نشان‌دهنده ماتریس ماسک توجه (علّی یا ماسک پرکردگی) است.
    
    \item \textbf{شاخه توجه خطی هسته‌محور (\lr{Linear Branch}):}
    با بهره‌گیری از نگاشت ویژگی فعال‌سازی $\phi(\mathbf{x}) = \operatorname{ELU}(\mathbf{x}) + 1.0 \ge 0$، ضرب ماتریسی بر اساس خاصیت شرکت‌پذیری بدون محاسبه صریح ماتریس $N \times N$ به شکل زیر محاسبه می‌شود:
    \begin{equation}
    \mathbf{O}_h^{\text{linear}} = \frac{\phi(\mathbf{Q}_h) \left( \phi(\mathbf{K}_h)^T \mathbf{V}_h \right)}{\phi(\mathbf{Q}_h) \left( \sum_{i=1}^N \phi(\mathbf{K}_{h, i}) \right)^T + \epsilon}
    \end{equation}
    در مدل‌های خودبازگشتی (\lr{Autoregressive Decoders})، این شاخه به صورت تکه‌ای (\lr{Chunked causal linear scan}) با انباشت حالات مخفی $\mathbf{S}_k = \mathbf{S}_{k-1} + \phi(\mathbf{K}_k)^T \mathbf{V}_k$ در ابعاد زمانی به صورت خطی پیاده‌سازی می‌گردد.
\end{enumerate}

گیت تلفیق سرها به ازای بردار ویژگی هر موقعیت $\mathbf{x}_t$ به صورت زیر است:
\begin{equation}
g_h(\mathbf{x}_t) = \sigma\left(\mathbf{w}_h^T \mathbf{x}_t + b_h\right), \quad b_h \approx 10.0
\end{equation}
خروجی نهایی پیش از تصویرسازی نگاشت پایانی از برهم‌نهی وزندار این دو شاخه حاصل می‌شود:
\begin{equation}
\mathbf{O}_h = g_h \odot \mathbf{O}_h^{\text{softmax}} + (1 - g_h) \odot \mathbf{O}_h^{\text{linear}}
\end{equation}
\begin{equation}
\mathbf{Y} = \operatorname{Concat}(\mathbf{O}_1, \dots, \mathbf{O}_H) \mathbf{W}^O
\end{equation}


\paragraph{تحلیل طیفی و نامساوی‌های پایداری خطا}
برای بررسی میزان دقت تقریب ماتریس وزن یا نگاشت توجه در ماژول \lr{hybrid\_attention}، تفاضل ماتریس فاکتورگیری‌شده از نگاشت متراکم اصلی بر اساس نامساوی هافدینگ-فون‌نیومان ارزیابی می‌گردد. فرض کنید $\mathbf{W} \in \mathbb{R}^{d_{\text{in}} \times d_{\text{out}}}$ ماتریس وزن هدف باشد که توسط تقریب کم‌مرتبه یا پراکنده با رتبه موثر $r$ جایگزین شده است. کران بالای خطای بازسازی بر مبنای نرم فروبنیوس به صورت زیر تصریح می‌شود:
\begin{equation}
\label{eq:spectral_bound_hybrid_attention}
\| \mathbf{W} - \mathbf{W}_{r} \|_F = \sqrt{\sum_{i=r+1}^{\min(d_{\text{in}}, d_{\text{out}})} \sigma_i^2(\mathbf{W})} \le \epsilon \| \mathbf{W} \|_F
\end{equation}
که در آن $\sigma_i(\mathbf{W})$ بیانگر $i$-امین مقدار منفرد مرتب‌شده ماتریس به صورت نزولی است. نسبت حفظ انرژی طیفی $\eta_r = \frac{\sum_{i=1}^r \sigma_i^2}{\sum_{i=1}^{d} \sigma_i^2}$ تضمین می‌کند که با انتخاب آستانه $\tau \ge 0.90$، بیش از نود درصد از پراکندگی واریانس داده‌ها حفظ شده و افت سیگنال در انتشار گرادیان رو به عقب به حداقل ممکن کاهش می‌یابد.

\paragraph{تحلیل مجانبی و پیچیدگی اشتراک حافظه سخت‌افزاری}
از منظر پیچیدگی محاسباتی و انتقال داده‌ها در سلسومراتب حافظه (حافظه سراسری \lr{HBM} به حافظه محلی \lr{SRAM})، عملیات ضرب ماتریسی در این مکانیزم نیازمند تخصیص بافر موقت با ابعاد فاکتورگیری‌شده است. برای پردازش دنباله‌ای به طول $N$ با بعد مخفی $D$، بار محاسباتی هسته از فرم درجه دو به رابطه زیر تقلیل می‌یابد:
\begin{equation}
\label{eq:complexity_hybrid_attention}
\mathcal{C}_{\text{hybrid\_attention}} = \mathcal{O}\left( B \cdot N \cdot D \cdot r + B \cdot H \cdot N \cdot d_k \right) \ll \mathcal{O}\left( B \cdot H \cdot N^2 \cdot d_k \right)
\end{equation}
این بهبود مرتبه پیچیدگی، با کاهش چشمگیر حجم تبادل تانسورها در گذرگاه‌های ارتباطی کارت گرافیک، مانع از اشباع پهنای باند حافظه (پدیده \lr{Memory-Bound}) در فرآیند استنتاج طولانی‌مدت می‌گردد.
